%% kapa: exam=1 round=32 year=2021 period=1 form=A subject=civil_law subject_ko=민법 seq=30 number=30
%% source: Q-Net 감정평가사 기출문제 (cst00309) artlSeq=5208397 file=extracted/1008/5208397/A형/2021년도 제32회 감정평가사 1차 1교시 A형.pdf
%% conversion: kordoc 4.24.51 -> tools/gen1.py (auto); review=proofread-sonnet; flags=none
\begin{question}{30}
공유물분할에 관한 설명으로 옳지 않은 것은? (다툼이 있으면 판례에 따름)
\begin{choices}
  \item 공유물분할의 효과는 원칙적으로 소급하지 않는다.
  \item 재판에 의한 공유물분할은 현물분할이 원칙이다.
  \item 공유관계가 존속하는 한, 공유물분할청구권만이 독립하여 시효로 소멸하지는 않는다.
  \item 공유토지를 현물분할하는 경우에 반드시 공유지분의 비율대로 토지 면적을 분할해야 하는 것은 아니다.
  \item 공유물분할의 조정절차에서 공유자 사이에 현물분할의 협의가 성립하여 조정조서가 작성된 때에는 그 즉시 공유관계가 소멸한다.
\end{choices}
\end{question}
\answer{5}
